\chapter[Mass Media Create a Shared Reality]{How Mass Media Create a Shared Reality}
\section{A Three-Day-Old Newspaper}
First, pick up a newspaper. It need not be new: find any paper from three days ago.

You immediately notice its feel. The paper is thin and slightly brittle, fuzzy at the folds, with ink faintly staining your fingertips. Notice something more interesting: its value is visibly falling. Three days ago, it was worth reaching into your pocket for money. Today, it is worth roughly as much as paper for wrapping fruit. Not one word has changed; what changed is its distance from now. News is a commodity with an extraordinarily short shelf life, measured in hours.

Now unfold it and inspect the front page. Suppose it has three stories: a distant earthquake, a newly passed law, and a celebrity's marital breakup. Ask yourself: why these three? Hundreds of thousands of things happened in this city yesterday. Someone was born, someone lost a job, people argued on a corner, a factory secretly discharged wastewater. Why did these three receive the most prominent positions while hundreds of thousands of others were treated as though they never happened?

A stranger question follows. Yesterday morning, millions of people in this city learned almost simultaneously about the same three things and assumed others knew too. Entering your classroom, you can discuss the earthquake with your deskmate without first checking whether they heard about it. You are certain they did. That certainty is so ordinary that you overlook its miraculousness: before newspapers, human history had no device allowing millions to share the same daily account of what happened in the world today.

This brittle sheet accomplished something astonishing: it installed a shared reality in an entire city, even a continent. This chapter concerns the machine---newspapers, radio, television, the internet---that manufactures the world everyone knows inside our heads, who controls it, and what can go wrong.

\section{A World for a Penny}
Enter a scene with me.

Time: an early morning in September 1833. Place: a New York street near Printing House Square. Dawn is breaking; horse manure and river water scent the air, and dockworkers' unloading calls sound from the waterfront.

At a corner stands a thinly dressed newsboy in his early teens, clutching newspapers still warm from the press and shouting, ``The Sun! Only a penny! The New York Sun!''

Nearby, another newspaper's sales counter is deserted. Its papers cost six cents, no small amount for an ordinary worker then, roughly a meal's price. Its readers are mainly merchants, lawyers, and politicians, so its content revolves around shipping schedules, stock prices, and congressional debates. A newspaper in that age is a class credential: which one you can afford reveals which world you inhabit.

The Sun costs a penny. Its publisher, Benjamin Day, calculates differently. Others earn money selling papers; he sells them so cheaply that they earn almost nothing, attracts thousands of readers, then sells their eyes to advertisers. The arrangement sounds alarmingly familiar today: almost every free phone app uses the same ledger. But in 1833, it was an invention.

Earning that penny also required different content. Congressional debates were too long and dull for ordinary readers. Instead came police-station fights, unusual court cases, street fires, love, and fraud. Two years later, the paper even published a series claiming astronomers had discovered winged creatures on the moon, complete with convincing details and citations to a nonexistent academic journal. New Yorkers rushed to buy it, driving circulation to the world's highest. When the fabrication emerged, readers were not especially angry: a penny's entertainment hardly seemed a loss.

Remember this corner. Three linked developments occur here: newspapers become affordable to everyone, content follows ordinary readers' interests rather than elite judgments about appropriate topics, and publishers discover their real commodity is attention, not news printed on paper. None has disappeared. All merely changed clothes and survived into today's phone notifications.

\section{One World, Different Versions}
Now state the problem without answering yet.

Newspapers' spread brought something apparently wonderful: information ceased to be the rich person's privilege. Dockworkers and bankers read the same front page; rural and urban people discussed the same election. A vast shared information space emerged, with common facts, topics, and emotions.

Place two different papers side by side, however, and cracks appear. The same strike becomes rioters disrupting order in a business-friendly paper and police suppressing petitioners in a labor-friendly one. The same army advances heroically in its home country's papers and flees in panic in its enemy's. Neither paper lies from beginning to end: some strikers did throw stones, and police did open fire. But each selects part of the truth, cuts away another, and assembles pictures pointing opposite ways.

A genuine question emerges:

\textbf{If everyone sees only part of the truth, where is our shared reality?}

Further, what role belongs to whoever decides which truth deserves the front page: editor, owner, censor, or today's recommendation algorithm? A neutral courier delivering events unchanged, or an invisible legislator deciding daily what matters, what is normal, and what is not worth knowing?

If the former, we can relax and treat media as a clean window. If the latter, the situation is serious: we think we see the world through glass, but have always been looking at pictures someone else hung for us.

This concerns more than news. It concerns elections' influence, wars' initiation, prejudices' reproduction across generations, and something close to you: whether that notification wants you to know something or become something.

Before continuing, answer privately: are media mainly windows or picture frames?

\section{Mirror or Mold?}
Two broad positions deserve their fullest presentation.

\textbf{Position one: media are a mirror, responsible only for reflecting the world as it is.}

Make this case fully. Intellectuals often mock it too quickly, unfairly.

Its first argument is that markets punish distortion. A paper that repeatedly publishes falsehoods loses readers. One warped reflection is a mistake; after ten, nobody looks. Competition corrects errors: ten papers watching each other prevent prolonged deception. Second, readers are not dough. Ordinary people have life experience. If a newspaper says the sky has fallen, they step outside, see it still there, and distrust it next time. Imagining people as empty vessels believing anything media say is itself elite arrogance. Historically, that arrogance often opens censorship's door: if people are so easily deceived, we must filter information for them. Third comes a painful question: if media are molds rather than mirrors, who should shape the mold---government, experts? Correcting media almost always becomes power taking them over. Of two evils, an occasionally distorted mirror is preferable to a mold controlled by others.

This powerful case is essentially the philosophical foundation of modern press freedom.

\textbf{Position two: media have never been mirrors. They are molds, and pretending to be mirrors is their greatest power.}

This case is equally complete. First, selection itself shapes. Mirrors need not choose what to reflect; media must choose. Countless daily events compete for headlines that accommodate fewer than one in a thousand. Unwritten selection standards---what counts as news, what deserves anger---shape readers more deeply than any individual story. Second, readers' ability to correct errors has a fatal gap. They can check matters directly connected to their lives: whether the sky fell is visible. But a distant war, financial policy, or country they will never visit offers no independent checking channel; media become their entire reality. Precisely where matters are most important, the mirror is hardest to test. Third, media's strongest influence is not persuading you of falsehoods, but keeping certain things absent from your mind altogether. The mirror account explains false news, not missing news. A city's newspapers might ignore a neighborhood for a century, effectively erasing its residents from the political map. A mold works most quietly through omission.

Elsewhere in the world, the same difficulty wears local faces.

In the years following Japan's Meiji Restoration, newspapers proliferated. Government subsidized supportive papers while legislating against critical ones. Readers on Tokyo streets could enter entirely different opinion worlds by buying different newspapers about the same political dispute. After Indian independence, state-run All India Radio reached across the country, bringing a single voice simultaneously into villages amid thousands of dialects. It spread agricultural and public-health knowledge, but opposition parties complained for decades that the governing party largely determined pre-election broadcasting. In Latin America, Mexican writer Miguel Sabido discovered in the 1970s and 1980s that television serials could quietly become textbooks. His popular dramas sent characters to adult literacy classes, followed by marked increases in real enrollments. Molds can certainly produce good things, but molding good things and manufacturing falsehoods use the same machine.

The darkest evidence comes from radio's golden age. In Rwanda in 1994, a station repeatedly called Tutsi people cockroaches, packaging hatred in popular songs and casual banter. Before killing began, the machine had performed the killers' hardest task: redefining neighbors as pests. Remember not merely that radio is frightening, but the mechanism. Years of ordinary, almost unobtrusive redefinition preceded catastrophe. Shared reality was rewritten first; reality then collapsed. Victims were not destroyed by one sentence. A million ordinary programs had first quietly revoked their humanity.

Both sides therefore hold half the truth. The mirror side rightly says public argument and competition correct many errors, and entrusting power to correctors is extremely dangerous. The mold side rightly says no mirror lacks an angle, and that angle is most dangerous when it claims not to exist.

\section[Three Questions about News]{Thinking Bridge: Take News Apart with Three Questions}
After positions comes a portable tool. Communication studies, the study of information flowing among people, has refined three concepts over a century. Do not fear their names: they are simply three questions to ask about any news item, trending topic, or notification.

\textbf{Question one: what does it want me to think about? Agenda-setting.}

Researchers found that media may struggle to tell you how to think but excel at telling you what to think about. The test is simple: follow headlines for a week, then compare conversations on streets, at dinner tables, and in class chats. The overlap is striking. A topic's weight on the page becomes its weight in public minds. A familiar paraphrase says newspapers are astonishingly successful at telling people what to think about, though often unsuccessful at telling them how to think. This is \textbf{agenda-setting}: a giant hand orders society's daily list of concerns. Listed topics draw everyone out; unlisted ones might as well never have happened.

\textbf{Question two: how does it want me to think? Framing.}

One event fits entirely different narrative frames. Factory layoffs can be presented as a troubled company's painful survival measure or as workers abandoned without livelihoods. The facts remain identical, but the frame assigns responsibility, directs emotion, and identifies blame. \textbf{Framing} is news's viewfinder: what enters, what is cropped, who occupies the center, whose voice tells it. Try a simple exercise: if someone with another position told this story, which sentence would move to the beginning? Whatever can be repositioned has been framed.

\textbf{Question three: who decides what I can see? Gatekeeping.}

Information passes countless doors between an event and your eyes: a reporter decides whether to investigate, an editor whether to publish, an editor-in-chief whether to release it, a platform algorithm whether to recommend it, and censorship whether it may remain. At each door stands a \textbf{gatekeeper}. The term came from a scholar of group psychology who first studied whose permission food requires between garden and table, then found information works similarly. Gatekeeping need not be conspiratorial; usually it mixes professional habits, business calculations, and deadlines. But not a conspiracy does not mean without consequences.

Together these questions form a sharp instrument for dissecting news. \textbf{Agenda-setting} concerns what to think about; \textbf{framing}, how to think; \textbf{gatekeeping}, who may speak and be heard. Next time a topic sweeps your class or news starts a family argument, ask all three. Many disputes reveal two sides fed different agendas and enclosed in different frames, each convinced it faces the event itself.

\section{A Public Relations Founder's Candid Account}
Now read a short primary source. Its rarity lies in the speaker: not a critic of media, but an acknowledged founder of the trade, speaking with remarkable frankness.

Edward Bernays, born in Austria-Hungary and later settled in America, pioneered public relations and was Sigmund Freud's nephew. In 1928 he published Propaganda. Its opening says, in paraphrase:

\begin{quote}
Deliberate, skillful manipulation of the public's organized habits and opinions is an important element of democratic society. Those manipulating this hidden social mechanism form an invisible government, the true governing power of our country. ---Edward Bernays, Propaganda (1928), Chapter 1
\end{quote}
Notice his tone. He is neither confessing nor exposing wrongdoing; he is \textbf{advertising his services}. His achievements included challenging the taboo against women smoking publicly for a tobacco company. He paid fashionable women to light cigarettes as torches of freedom in an Easter parade and alerted photographers. For a bacon producer, he promoted a hearty American breakfast as a health standard by securing thousands of doctors' questionnaire signatures. He rarely fabricated facts outright. His higher-level craft arranged scenes, borrowed authority, and created images so the public believed it had reached conclusions independently.

Reading this century-old account, do two things. First, notice invisible government. Bernays thought democracy too large and complex for everyone to decide everything directly. Professionals therefore had to organize public attention backstage. He considered this democracy's lubricant, even a virtue. Second, weigh the danger: if manipulating the public is democracy's necessary component, is a voter's ballot the product of personal thought or organized influence? Bernays was comfortable with the contradiction. You need not be.

Incidentally, the book later had a notorious reader: Nazi Germany's propaganda minister reportedly studied and borrowed Bernays's techniques. The same tools can sell cigarettes, promote public health, or prepare hatred's path. Tools do not choose their owners, a fact we will repeatedly encounter.

\section{Business, Power, or the Machine?}
After Bernays, you probably accept that media shape reality. The next question is \textbf{why}. More than two competing explanations address the same fact, each with strengths and blind spots.

\textbf{Explanation one: business. Media are commodities, and attention is what they truly sell.}

This explanation stretches directly from that penny-paper corner. Newspapers must earn money, money requires readers, and papers publish what readers enjoy. Once advertising dominates revenue, the logic becomes packaging as many eyes as possible for advertisers. Distortion here arises not from malice but market preference: sensation outsells stability, conflict outsells reconciliation, simple stories outsell complex reality. The old saying that a dog biting a man is not news but a man biting a dog is captures news's preference for the abnormal. Daily readers therefore systematically overestimate danger and disorder. This powerfully explains a century of sensationalism, clickbait, and catastrophe-driven reporting. Its limit is explaining things said \textbf{even without profit}. Why would papers lose money to attack a party? Why do some countries regard loss-making media as a normal public expense? Business logic reaches its boundary here.

\textbf{Explanation two: power. Media are instruments of rule, serving their owners or controllers.}

Rather than trace content, trace ownership: whose newspaper, whose television station, whose instructions a platform follows. Rulers need shared reality to maintain order. Ancient Chinese dibao, political bulletins copied within bureaucracy whose readership eligibility was itself power, belong to the same family as wartime censorship and Bernays's invisible government. This sharp account explains what business cannot: why obviously newsworthy matters vanish collectively, or differently aligned outlets suddenly speak alike on particular issues. Its limit is sliding into an everything-is-manipulation conspiracy, erasing journalistic dignity, media exposing one another, and readers' judgment. In reality, power's desire to control and markets' desire to sell frequently clash; neither wins everything.

\textbf{Explanation three: the machine. Regardless of users or purposes, the medium's technical form transforms content.}

This is the most abstract and surprising explanation: change the channel rather than the content, and the world still changes. One speech printed on paper, spoken by a radio announcer, or cut into a fifteen-second video becomes three experiences. With every word unchanged, audiences' understanding, emotion, and memory have altered. Here business and power merely ride technology; the silent director is the medium's form. This explains what the others struggle with: why television suddenly made politicians' looks, speaking skills, and camera presence as important as policies worldwide. Voters did not simply grow shallow; cameras receive images and emotions, not dissertations. Yet technological explanation can make people too passive, as though form predetermines everything and institutions, money, and human choices count for nothing. That plainly goes too far: the same internet develops entirely different things in different societies.

Each explanation illuminates something and misses something. A mature view does not select just one. It asks which forces operate behind this particular story: who earned money, who gained power, and what the machine quietly changed.

\section[The Medium Is the Message]{Deeper Challenge: The Medium Is the Message}
The third explanation deserves deeper examination. It comes from twentieth-century communication studies' most famous and frequently misread slogan: the medium is the message.

Canadian scholar Marshall McLuhan proposed it. He did not mean content is unimportant. Rather, \textbf{a new medium changes the world most through how it exists, not what it carries}. Consider his frequent example, electric light. It transports no content; it simply shines at night. Yet it enables night shifts and evening sports, rewriting cities' timetables. This contentless medium changes social structure more deeply than many content-bearing ones. Through this lens, media history changes protagonists. Instead of which news story influenced history, ask how printing, telegraphy, radio, television, and the internet each reshaped minds and society.

A field called media ecology pursues this line. Its founders observed that in purely oral societies, knowledge survives through rhythm, proverbs, and repetition. Homeric epics can be sung, not looked up; elders are libraries and memory is culture's lifeline. Writing lets knowledge leave people and wait on paper for checking. Concepts of author and original text emerge; laws become written clauses and logic can be repeatedly examined. The cost is solitary, silent reading displacing communal chanting. Printing drastically lowers textual costs. Within a century of the Gutenberg Bible, Europe develops people who read scripture themselves, publish newspapers themselves, and believe in their own judgment. Historians widely consider printing essential to the Reformation and modern science, not simply because it spread particular views, but because books in everyone's hands rearranged authority. Then telegraphy makes information travel faster than people for the first time; distant events compete with nearby ones for attention. News as a business is, essentially, telegraphy's child.

With television, Neil Postman voiced a sharp concern. In paraphrase, print trains people to read long sentences, follow arguments, and tolerate dullness. Television turns everything into fast food of images and emotion. When politics, religion, and education must wear entertainment's clothes to be seen, a format quietly degrades society's capacity for serious discussion. His target is not particular bad programs, but a television \textbf{form} welcoming only certain content. This is a forceful application of the medium-is-the-message argument.

The theory has faults. Overextended, it becomes technological determinism, treating humans as media's puppets. As noted, identical technologies bear different fruit in different societies. Yet it teaches an irreplaceable question. Before arguing about a new medium's content, ask: \textbf{what does this form favor or exclude, how does it divide time, and what price does it place on attention?}

Apply that question to today's internet and recommendation algorithms and you immediately feel its edge. Newspaper editors gatekeep through professional judgments that are at least stable, public, and open to criticism. Algorithms gatekeep through your past clicks and viewing time, giving everyone a different front page without explanation, signature, or office hours. You cannot even find an editor-in-chief to send an angry letter. McLuhan did not reach the algorithmic age, but his slogan sounds like its advance warning: you think you choose among content, while form has already made the choice.

\textbf{An important counterexample belongs here, because limitless media power is among this field's easiest misjudgments.}

On the eve of Halloween in 1938, CBS broadcast The War of the Worlds, a radio drama adapted from science fiction that imitated news bulletins reporting a Martian invasion. Newspapers erupted the next day: they described nationwide panic, fleeing crowds, and people nearly driven to suicide by fear. For decades, the story served as proof of radio's ability to manipulate the masses; you may have encountered it. Later communication historians revisited archives and found a far less frightening picture. Fewer people listened than legend claimed; truly panicked listeners were fewer still. Most knew it was drama or verified it by telephone. \textbf{Newspapers} inflated much of the panic afterward, while viewing radio as a deadly rival for advertising. Depicting it as an irresponsible deception machine materially benefited the press.

This nearly perfect counterexample teaches two lessons. First, it corrects the facts: people did not collectively lose their minds upon hearing radio; they had checking abilities, social networks, and skepticism. Second, it enacts the chapter's theme: the story that media manufactured panic was itself a media-manufactured shared reality, persisting for decades because it served one side's interests. Even our sense of how frightening media are comes from media. Remember this double structure. Overestimating their power is as mistaken as underestimating it, and behind either mistake often stands a benefiting gatekeeper.

\section{Are Our Feeds Still the Same World?}
Now return to today, and your pocket.

Your phone's world is related to every machine discussed here. Notifications resemble newsboys' cries; trending lists are millions' common front page; short videos push television's fast-food tendency to its limit; algorithms fulfill Bernays's dream of an invisible government that never sleeps, takes no salary, and customizes reality for each person. But one change is unprecedented: \textbf{the shared information space is fragmenting}.

Newspaper-era distortion gave a whole city the same distortion; at least everyone argued about the same thing. Algorithmic distortion gives each person a custom reality. Grandpa's phone fills with health rumors and international conspiracies, Mom's with educational anxiety, and yours with games, celebrities, and school gossip. Three people at one dinner table lift their heads from three almost nonintersecting worlds. Past family arguments over news at least proved a shared information sky. Today, families may lack even common material for argument: nobody else has heard of the event making you furious.

School works similarly. A rumor about a student circulates hundreds of times in a group within half a day, faster than any nineteenth-century newspaper could imagine. Its correction receives a fraction of the views. This is not an individual moral defect but a mechanism: rumors offer conflict, emotion, and suspense, fitting the ancient criteria of what sells. Platforms feed them to more people, who enlarge them further. Agenda-setting, framing, and gatekeeping remain alive, but identifiable editors have become recommendation systems nobody can name.

Another player is AI. At this book's writing, machines can mass-produce convincing text, images, and sounds. Manufacturing a shared reality once required buying presses or founding stations; now the cost approaches zero. As pictures cease to guarantee truth, the photographs, recordings, and videos used to check reality for over a century rapidly lose value. Rwanda already answered why shared information matters: it underpins strangers' cooperation, voting, and reconciliation. Your phone daily answers why it is fragile: this foundation never grew naturally. It consists of costly, fragile institutions requiring continual maintenance---honest reporting habits, accountable gatekeepers, willing fact-checking readers. Maintenance benefits everyone; destruction can profit individuals.

But portraying today as entirely dark would also distort it. The same internet connects rare-disease patients, brings ignored towns to national attention, and grants ordinary people distribution power once reserved for editors-in-chief. You are a potential gatekeeper. The mold account's deepest insight becomes hope: if angle-free news does not exist, the solution is not to find angle-free media but to make as many angles as possible coexist visibly and check one another. A broken mirror is not the danger. The danger is one remaining shard claiming to be the entire sky.

\section{Who Gets the Key?}
We have traveled far from a penny-paper corner to your phone, but one opening must remain unclosed.

Shared reality requires gatekeepers. Information is too abundant and time too short; someone or something must filter for us. A world without gatekeepers is not freedom but a flood. Yet gatekeeping inherently carries power: deciding what hundreds of millions consider, resent, or forget.

One question remains, with no standard answer:

\textbf{If someone must hold the filtering key---editors, government, platform algorithms, or ordinary users like you---whom would you choose? What rules should they follow? And who should oversee this overseer, by what means?}

Give reasons, and offer one plausible reason for the option you oppose most. If every option seems fatally flawed, congratulations: you have not answered incorrectly, but touched the problem's true shape. Humanity has struggled with it for nearly two hundred years. The best response so far is not a perfect answer, but continued argument and refusal to let anyone claim to have found that answer for you.
